% ===================================================================
%  ТАБЛИЦА № 7 — Деревня зимой. — Ферма весной.
%  Traduction 2026 d'apres texto/io, controlee sur texto/fr, texto/en
%  et texto/es. Consigne : voir texto/ru/10-tablica-01.tex.
% ===================================================================

%%K t07-noto-1 noto
\VUnotes{143.2mm}{%
(*) Подробности обеда см. в таблицах 6 и 13.%
}

%%K t07-apar-1 apar
\VUsaut{4.0mm}
\VUcentre{13.2pt}{90}{ВТОРАЯ СЕРИЯ}
\VUsaut{3.0mm}
\VUfilet{16mm}
\VUsaut{5.6mm}
\VUtitre{15.0pt}{60}{ТАБЛИЦА № 7}
\VUsaut{3.0mm}

%%K t07-tit-1 sub
\VUcentre{12.0pt}{0}{\textit{Первая сцена.}}
\VUsaut{3.0mm}

%%K t07-tit-2 sub
\VUpk{10.2pt}{40}{Деревня зимой.}
\VUsaut{4.0mm}

%%K t07-c1-01-1 p
1. --- На новогодних праздниках я поехал с отцом в Новосёлки, деревушку,
где у него были дела.

%%K t07-c1-02-1 p
2. --- Мы поехали в \VUgras{дилижансе} \textsuperscript{(11)}, который ходит между
городом и этой деревней; но так как было очень холодно, поездка в столь
неудобной повозке вышла совсем не приятной.

%%K t07-c1-03-1 p
3. --- Так что мы искренне обрадовались, когда \VUgras{кучер} остановил
свой дилижанс на площади, перед \VUgras{трактиром} \textsuperscript{(2)}
\textit{«Белый Медведь»}. \VUgras{Трактирщик} \textsuperscript{(4)} поджидал нас
на пороге, спокойно куря свою \VUgras{трубку}.

%%K t07-c1-03-2 p
Он пригласил нас войти, и меня не пришлось просить дважды, чтобы я уселся
перед огнём из виноградной лозы, трещавшим в очаге. Тогда мне стал
нипочём резкий \VUgras{северный ветер}, от которого скрипела старая
\VUgras{вывеска} \textsuperscript{(3)} и который уносил чёрными клубами
\VUgras{дым} \textsuperscript{(1)}, выходивший из трубы.

%%K t07-c1-04-1 p
4. --- Был час обеда. Не мешкая, мы отдали должное простой, но вкусной
еде, которую нам подали (*).

%%K t07-tit-3 sub
\VUpk{10.2pt}{40}{Несколько визитов.}
\VUsaut{3.0mm}

%%K t07-c1-05-1 p
5. --- После обеда я пошёл с отцом, страховым агентом, к его клиентам.
Сперва мы посетили \VUgras{кузнеца} \textsuperscript{(5)}, который живёт рядом с
трактиром. Он подковывал лошадь. Я подивился силе его подручного, который
крепко держал \VUgras{копыто} лошади у своего кожаного фартука, пока
кузнец тяжёлыми ударами молота прибивал \VUgras{подкову} \textsuperscript{(6)}.

%%K t07-c1-06-1 p
6. --- Потом мы пошли к деревенскому \VUgras{сапожнику} \textsuperscript{(7)},
старому Якову. Хотя вывеской ему служил великолепный \VUgras{сапог}, он
довольствовался тем, что подшивал старую пару стоптанных башмаков.

%%K t07-c1-06-2 p
Старик со смехом сказал нам: «В городе я был „сапожником“, и у меня не
было заказчиков. Здесь я „починщик“, и работы хоть отбавляй».

%%K t07-c1-07-1 p
7. --- Он говорил нам о своём соседе \VUgras{цирюльнике} \textsuperscript{(8)},
клиенты которого все недовольны. Его \VUgras{бритвы} всегда с зазубринами,
а \VUgras{ножницы} никогда толком не режут. Но так как он единственный
цирюльник в деревне, людям, понятно, ничего не остаётся, как бриться и
стричься у него. К тому же он человек скупой и неприветливый.

%%K t07-c1-08-1 p
8. --- К счастью, другие \VUgras{соседи} на него не похожи.
\VUgras{Колёсник} \textsuperscript{(9)}, например, человек добрый, работящий и
услужливый. Отдать ему телегу в починку --- одно удовольствие. Работа у
него всегда чисто сделана.

%%K t07-c1-09-1 p
9. --- Не жаловался старый Яков и на \VUgras{шорника} \textsuperscript{(10)},
которому иногда помогает шить \VUgras{сбрую}, когда работа не терпит.

%%K t07-c1-09-2 p
Отец спросил его о семье. «Все здоровы. Внучка моя в \VUgras{школе}
\textsuperscript{(12)}. \VUgras{Учительница} \textsuperscript{(13)} ею очень довольна;
она хорошая \VUgras{ученица} \textsuperscript{(14)}. Не то что этот шалопай, мой
сын; у него только игры на уме. \VUgras{Учитель} \textsuperscript{(15)} никак не
заставит его выучить хоть что-нибудь».

%%K t07-tit-4 sub
\VUsaut{3.0mm}
\VUpk{10.2pt}{40}{Деревенские толки.}
\VUsaut{3.0mm}

%%K t07-c1-10-1 p
10. --- Затем старик рассказал нам деревенские новости. Собирались строить
новую школу, потому что та, что в \VUgras{сельской управе}, стала уж
слишком тесна. Это было бы нетрудно, ведь довольно было купить часть сада
\VUgras{мельника}. Бедняге это было бы очень кстати. С самых холодов
\VUgras{жернова} \VUgras{мельницы} \textsuperscript{(16)} не могли молоть зерно,
потому что \VUgras{колесо} \textsuperscript{(17)} схватило \VUgras{льдом}
\textsuperscript{(25)} на \VUgras{ручье} \textsuperscript{(23)}.

%%K t07-c1-11-1 p
11. --- «Кстати о постройках: видели вы нашу новую \VUgras{церковь}
\textsuperscript{(18)}? Она очень живописна под снежной шапкой. \VUgras{Вороны}
\textsuperscript{(19)}, которые не узнают больше своей старой колокольни, уныло
садятся на большое дерево \VUgras{кладбища} \textsuperscript{(20)}».

%%K t07-c1-12-1 p
12. --- Упоминание о кладбище напомнило сапожнику о тех знакомых отца,
которые умерли за минувший год. «Сколько \VUgras{могил} \textsuperscript{(21)},
сударь мой, прибавилось к прежним под \VUgras{кипарисом} \textsuperscript{(22)}!
Смерть взяла нашего старого нотариуса. Вся деревня была на его
\VUgras{похоронах}».

%%K t07-c1-13-1 p
13. --- «Вон его преемник катается на ручье. Он давеча заходил, чтобы я
починил ремешок его \VUgras{конька} \textsuperscript{(26)}, который лопнул».

%%K t07-c1-14-1 p
14. --- «А вон там, на берегу ручья, новый \VUgras{сельский сторож}
\textsuperscript{(27)}. Он бывший жандарм и гроза деревенских шалопаев. Они
разбегаются, едва завидят его \VUgras{гетры} \textsuperscript{(29)} и его
\VUgras{фуражку} \textsuperscript{(28)}».

%%K t07-c1-15-1 p
15. --- «А! вон старый Иосиф везёт \VUgras{воз хвороста} \textsuperscript{(30)}
булочнику. --- И верно, сказал отец, узнаю его упряжку \VUgras{мулов}
\textsuperscript{(31)}. Мне нужно с ним переговорить, и я этим воспользуюсь.
Прощай, старый Яков».

%%K t07-c1-16-1 p
16. --- Как раз когда мы выходили, деревенские дети выбегали из школы и
кидались на \VUgras{ледяную дорожку} \textsuperscript{(33)}, рискуя упасть и
сломать что-нибудь при \VUgras{падении} \textsuperscript{(34)}. Один из них
взял у колёсника свои \VUgras{санки} \textsuperscript{(32)}, посадил в них
товарища и побежал.

%%K t07-c1-17-1 p
17. --- Другие бросились к \VUgras{сугробу} \textsuperscript{(37)} возле
\VUgras{моста} \textsuperscript{(40)} и принялись обстреливать \VUgras{снежную
бабу} \textsuperscript{(39)}, которую слепили накануне и которой дали шляпу,
трубку и школьную сумку через плечо.

%%K t07-c1-18-1 p
18. --- Им был нипочём ледяной ветер и стужа. У большинства не было даже
\VUgras{шарфа} \textsuperscript{(35)}, а чтобы согреться после игры в
\VUgras{снежки} \textsuperscript{(38)}, им довольно было горсти обжигающих
\VUgras{каштанов} \textsuperscript{(42)}, купленных у старой Жаклины,
\VUgras{торговки}, которая дрожит от холода под своей слишком тонкой
\VUgras{шалью} \textsuperscript{(43)}.

%%K t07-tit-5 sub
\VUsaut{3.0mm}
\VUcentre{12.0pt}{0}{\textit{Вторая сцена.}}
\VUsaut{3.0mm}

%%K t07-tit-6 sub
\VUpk{10.2pt}{40}{Ферма весной.}
\VUsaut{4.0mm}

%%K t07-c2-01-1 p
1. --- При всех радостях зимы мы с глубоким весельем встречаем возвращение
\VUgras{весны}.

%%K t07-c2-01-2 p
Какая отрадная картина --- двор фермы вечером, в мае!

%%K t07-c2-02-1 p
2. --- На горизонте медленно садится \VUgras{солнце} \textsuperscript{(1)}, заливая
\VUgras{поля} \textsuperscript{(3)} багряными \VUgras{лучами} \textsuperscript{(2)}.
Работящий \VUgras{пахарь} \textsuperscript{(6)} спешит вспахать своё
\VUgras{поле} \textsuperscript{(4)}.

%%K t07-c2-02-2 p
Своим \VUgras{плугом} \textsuperscript{(5)}, который тянет крепкая \VUgras{лошадь}
\textsuperscript{(7)}, он проводит \VUgras{борозду} \textsuperscript{(8)}, в которую
\VUgras{сеятель} \textsuperscript{(9)} горстями разбросает \VUgras{семя}, что
даст золотые колосья.

%%K t07-c2-03-1 p
3. --- \VUgras{Косцы} \textsuperscript{(11)} со своими косами скосили луговую
траву, ворошильщики высушили \VUgras{сено} \textsuperscript{(10)}, переворачивая
его \VUgras{вилами} \textsuperscript{(12)} и граблями, и вот они возвращаются
домой.

%%K t07-c2-03-2 p
Одни идут впереди тяжёлого воза, влекомого волами, неся орудия на плече.
Другие, ленивее, удобно устроились на душистом сене.

%%K t07-c2-04-1 p
4. --- Проходя, они дружески здороваются с \VUgras{садовником}
\textsuperscript{(16)}, который трудится в \VUgras{саду} \textsuperscript{(13)},
подрезая \VUgras{плодовые деревья} и прививая \VUgras{груши}
\textsuperscript{(14)}. \VUgras{Сливы} \textsuperscript{(15)} в цвету, а на хорошо
удобренном \VUgras{огороде} \textsuperscript{(17)} овощи, капуста и салат уже
хороши.

%%K t07-c2-04-2 p
На низкой ограде красивый \VUgras{павлин} \textsuperscript{(18)} распускает
хвост, чтобы показать свои великолепные перья.

%%K t07-tit-7 sub
\VUpk{10.2pt}{40}{Двор фермы.}
\VUsaut{3.0mm}

%%K t07-c2-05-1 p
5. --- Двери \VUgras{конюшни} \textsuperscript{(19)} открыты, и видны ясли и
\VUgras{подстилка} \textsuperscript{(23)} \VUgras{кобылы} \textsuperscript{(21)},
которую \VUgras{конюх} \textsuperscript{(20)} только что распряг.
\VUgras{Жеребёнок} \textsuperscript{(22)}, такой уморительный на длинных ногах,
идёт за матерью, резвясь.

%%K t07-c2-06-1 p
6. --- У \VUgras{колодца} \textsuperscript{(29)} \VUgras{коровы} \textsuperscript{(25)}
идут пить из деревянного \VUgras{корыта} \textsuperscript{(26)}, служащего им
водопоем. \VUgras{Телёнок} \textsuperscript{(24)} пользуется случаем пососать, а
\VUgras{вол} \textsuperscript{(27)}, почуяв добрый запах корма, весело мычит и
гордо поднимает голову с могучими \VUgras{рогами} \textsuperscript{(28)}.

%%K t07-c2-07-1 p
7. --- Все за работой. \VUgras{Работница} \textsuperscript{(32)} держит
\VUgras{верёвку} \textsuperscript{(31)} колодца. Она черпает воду \VUgras{ведром}
\textsuperscript{(30)}, чтобы напоить скотину. Возле \VUgras{амбара}
\textsuperscript{(33)} человек сгребает рассыпанную солому, а перед дверью
\VUgras{фермы} \textsuperscript{(34)} старая \VUgras{пряха} \textsuperscript{(35)} со
своей прялкой и \VUgras{куделью} прядёт \VUgras{лён} или \VUgras{коноплю},
как в былые дни.

%%K t07-tit-8 sub
\VUsaut{3.0mm}
\VUpk{10.2pt}{40}{Фермер.}
\VUsaut{3.0mm}

%%K t07-c2-08-1 p
8. --- \VUgras{Фермер} \textsuperscript{(36)} управляет всей этой работой и даёт
своим людям указания. Он велит убрать под навес \VUgras{борону}
\textsuperscript{(40)} и \VUgras{кирку} \textsuperscript{(39)}, оставленные возле
\VUgras{будки} \textsuperscript{(38)} \VUgras{собаки} \textsuperscript{(37)}.

%%K t07-c2-09-1 p
9. --- Его окрики вспугивают \VUgras{голубя} \textsuperscript{(41)}, который во
весь дух улетает в \VUgras{голубятню} \textsuperscript{(42)}, и \VUgras{кур}
\textsuperscript{(43)}, которые спешат обратно в \VUgras{курятник}
\textsuperscript{(44)}.

%%K t07-c2-10-1 p
10. --- Перед \VUgras{овчарней} \textsuperscript{(48)} старый \VUgras{пастух}
\textsuperscript{(45)} пригоняет своё \VUgras{стадо} \textsuperscript{(49)}. Держа
свой \VUgras{посох} \textsuperscript{(46)}, который не расстаётся с ним, как и
выцветшая \VUgras{блуза} \textsuperscript{(47)}, он ведёт в загон
\VUgras{баранов} \textsuperscript{(50)}, \VUgras{овец} \textsuperscript{(53)},
\VUgras{ягнят} \textsuperscript{(54)}, \VUgras{коз} \textsuperscript{(51)} и
\VUgras{козлят} \textsuperscript{(52)}. Его верный \VUgras{пёс} \textsuperscript{(55)}
Аргус идёт последним и лаем подгоняет отставших.

%%K t07-c2-11-1 p
11. --- Весь этот шум и суета не тревожат спокойствия доброй
\VUgras{дойной коровы} \textsuperscript{(56)}, которая позвякивает своим
\VUgras{колокольчиком} \textsuperscript{(57)}, пока её доят перед дверью
\VUgras{коровника} \textsuperscript{(58)}.

%%K t07-c2-12-1 p
12. --- Она словно понимает, что должна дать молока, чтобы
\VUgras{фермерша} \textsuperscript{(60)} могла сбить \VUgras{масло} в
\VUgras{маслобойке} \textsuperscript{(61)} или сделать отличные \VUgras{сыры}
\textsuperscript{(64)}, которые обсыхают на доске, положенной поперёк
\VUgras{кадки} \textsuperscript{(63)}.

%%K t07-c2-13-1 p
13. --- Всю \VUgras{молочную} \textsuperscript{(59)} держит в чистоте
\VUgras{батрак} \textsuperscript{(65)}, который только что вымыл \VUgras{бидоны}
\textsuperscript{(62)} в большом ведре, которое несёт.

%%K t07-tit-9 sub
\VUsaut{3.0mm}
\VUpk{10.2pt}{40}{Птичий двор.}
\VUsaut{3.0mm}

%%K t07-c2-14-1 p
14. --- В своих тяжёлых деревянных башмаках он ходит довольно неуклюже и
потому распугивает \VUgras{индюка} \textsuperscript{(68)}, \VUgras{уток}
\textsuperscript{(66)}, \VUgras{селезней} \textsuperscript{(67)} и \VUgras{гусей}
\textsuperscript{(69)}, которые широко разевают \VUgras{клювы} \textsuperscript{(70)}
и поднимают тревожный гам.

%%K t07-c2-14-2 p
\VUgras{Петух} \textsuperscript{(71)}, более воинственный, поднимает свой красный
\VUgras{гребень} \textsuperscript{(72)}, встряхивает перьями \VUgras{хвоста}
\textsuperscript{(73)} и издаёт звонкое «кукареку».

%%K t07-c2-15-1 p
15. --- \VUgras{Наседка} \textsuperscript{(74)} копается в \VUgras{навозной куче}
\textsuperscript{(81)} со своими \VUgras{цыплятами} \textsuperscript{(75)}.
\VUgras{Поросёнок} \textsuperscript{(78)} бежит к матери, огромной
\VUgras{свинье} \textsuperscript{(77)}, которую видно возле хлева.

%%K t07-c2-16-1 p
16. --- В своих клетках \VUgras{кролики} \textsuperscript{(79)} жадно едят нежную
\VUgras{траву} \textsuperscript{(80)}, которую приносит им дочка фермера. Женя
очень любит своих кроликов и каждый день, взяв в курятнике свежие
\VUgras{яйца} \textsuperscript{(76)}, приходит навестить своих маленьких друзей.
\VUsaut{6.0mm}
\VUfilet{22mm}
